\documentclass{situationist-preprint}

\title{Zeno's Paradoxes and the Possibility of Change}
\msid{MET-21-b}
\msfamily{MET-21}
\msdiscipline{Metaphysics}
\mstype{Companion manuscript}
\mskeywords{Zeno; motion; change; dichotomy; Achilles; arrow; supertasks; at-at theory}

\begin{document}
\maketitle

\begin{abstract}
MET-21-a defends the possibility of change against Parmenides by exposing an equivocation between the existential and the predicative uses of \enquote{is}. Zeno's arguments against motion do not rest on that equivocation, and a defence of change must answer them separately. We take the four arguments Aristotle reports. The Dichotomy and the Achilles assume that infinitely many tasks cannot be completed in a finite time; the assumption is false when their durations have a finite sum, and the objection from supertasks does not apply to continuous motion. The Arrow infers from the arrow's being at rest at each instant that it is at rest throughout; the inference fails on the at-at conception of change that MET-21-a already adopts. The Stadium refutes only a discrete space and time with a minimal speed. Change, including motion, is possible.
\end{abstract}

\section{Arguments that do not turn on \enquote{is}}

Parmenides' argument, as MET-21-a reconstructs it, moves from \enquote{what changes comes from what is not} to \enquote{what changes comes from nothing}, and fails because no single reading of \enquote{is} makes both premises true. Zeno's arguments, as Aristotle reports them \citep[VI.9]{AristotlePhys}, are of another kind. They take for granted that there are bodies and places and ask whether a body can pass from one place to another. Their premises concern the structure of space and time, and an answer to Parmenides leaves them standing. We therefore answer them in their own terms, using one element of MET-21-a: Definition~2.2 of that manuscript, on which an object changes with respect to $F$ if it is $F$ at one time and not $F$ at another.

\section{The Dichotomy and the Achilles}

To reach the end of a course, a runner must first reach the half-way point, and before that the quarter-way point, and so on; the course divides into infinitely many segments, each of which must be traversed, and infinitely many tasks cannot be completed. The Achilles has the same structure: the pursuer must first reach the point where the tortoise was, by which time the tortoise has moved on, and so infinitely often.

\begin{proposition}\label{p:sum}
If a runner moves at constant speed $v$ over a course of length $d$, the segments of the Dichotomy are traversed in times whose sum is $d/v$.
\end{proposition}

\begin{proof}
The $n$th segment has length $d/2^n$ and is traversed in time $d/(2^n v)$. The sum of these times over all $n \geq 1$ is $d/v$.
\end{proof}

The argument needs the premise that infinitely many tasks cannot be completed in a finite time. Proposition~\ref{p:sum} shows that the premise is false for the tasks in question, since their durations form a convergent series. The same holds of the Achilles, where the times taken to reach the tortoise's successive positions decrease geometrically and sum to the time at which the pursuer overtakes. \citet{Grunbaum1967} gave the modern form of this answer and showed that it requires only the standard measure-theoretic treatment of continuous intervals; the reply does not depend on any claim about what can be counted or surveyed.

Aristotle's own answer was different: the segments are potential, not actual, divisions, and a runner who does not stop at the midpoints does not perform infinitely many actual tasks \citep[VIII.8]{AristotlePhys}. The modern answer is stronger, since it allows that the divisions are real and shows that their number is no obstacle.

\begin{objection}[Supertasks]
\citet{Thomson1954} argued that infinitely many tasks in a finite time is incoherent: a lamp switched on and off at times $1/2$, $3/4$, $7/8$ and so on before a minute has elapsed is neither on nor off at the end. Motion through the Dichotomy is such a supertask.
\end{objection}

\begin{reply}
The lamp has no determinate final state because nothing in the sequence of switchings fixes one; its state at the end is not a limit of its states before. The runner's position, by contrast, is a continuous function of time, and its value at the end is the limit of its values at the midpoints. The incoherence Thomson found belongs to sequences whose final state is underdetermined by the series, not to supertasks as such, as \citet{Benacerraf1962} showed. Motion through the Dichotomy is completed because the runner's trajectory is continuous, and continuity is what fixes where she is when the series is over.
\end{reply}

\section{The Arrow}

At each instant of its flight, the arrow occupies a space equal to itself, and so, at that instant, it is not moving. Time is composed of instants. So the arrow is at rest throughout its flight.

The argument does not equivocate on \enquote{is}; it infers a property of an interval from properties of its instants. The inference is invalid on the conception of change that MET-21-a adopts. By Definition~2.2 of MET-21-a, change is being $F$ at one time and not $F$ at another; motion is being at one place at one time and at another place at another time. This is the at-at conception that \citet{Russell1903} defended: motion is not a state the arrow is in at an instant but a relation between its positions at different instants. Being at rest is likewise a matter of intervals: a body is at rest during an interval if it is at the same place at every instant of the interval. At a single instant, the arrow is neither moving nor at rest in these senses; it is simply somewhere. Its instantaneous velocity is defined, but as the limit of its average velocities over shrinking intervals, so it is a fact about the neighborhood of the instant and not about the instant alone.

Given the at-at conception, the premise that the arrow is at rest at each instant is either false, if \enquote{at rest} is read as a property of intervals, or true but irrelevant, if it means only that the arrow occupies a single place at that instant. On the second reading, the conclusion does not follow, since an object can occupy a single place at each instant and different places at different instants. The Arrow is the clearest case in which a theory of change, rather than a semantics of \enquote{is}, is what answers the Eleatic.

\section{The Stadium}

Three rows of bodies, one at rest and two moving in opposite directions at equal speeds, pass one another; a moving body passes twice as many bodies of the other moving row as of the stationary row in the same time, and Zeno concludes that half the time equals the double. On a continuous space and time, the argument confuses relative with absolute motion, and there is nothing paradoxical in it. On a discrete space and time, in which each body moves at most one place per indivisible instant, the argument shows that two bodies moving in opposite directions can pass each other without there being an instant at which they are aligned. That is a feature of discrete motion, and a reason to doubt that motion is discrete in this way, but not an argument that motion is impossible.

\section{Conclusion}

Zeno's arguments do not rest on the equivocation that undoes Parmenides, and they need separate answers. The Dichotomy and the Achilles assume that infinitely many tasks cannot be completed in a finite time, which is false for tasks whose durations converge; the supertask objection fails for continuous trajectories. The Arrow infers rest from instantaneous position, which the at-at conception of change blocks. The Stadium refutes at most a discrete space and time with a minimal speed. With these answers, the conclusion of MET-21 holds without qualification: change is possible.

\begin{thebibliography}{99}\small
\bibitem[Aristotle(Physics)]{AristotlePhys} Aristotle. \emph{Physics}, Books VI and VIII.
\bibitem[Benacerraf(1962)]{Benacerraf1962} Benacerraf, P. (1962). Tasks, super-tasks, and the modern Eleatics. \emph{The Journal of Philosophy}, 59(24), 765--784.
\bibitem[Gr\"unbaum(1967)]{Grunbaum1967} Gr\"unbaum, A. (1967). \emph{Modern Science and Zeno's Paradoxes}. Wesleyan University Press.
\bibitem[Russell(1903)]{Russell1903} Russell, B. (1903). \emph{The Principles of Mathematics}. Cambridge University Press.
\bibitem[Thomson(1954)]{Thomson1954} Thomson, J.~F. (1954). Tasks and super-tasks. \emph{Analysis}, 15(1), 1--13.
\end{thebibliography}

\end{document}
